\documentclass[sigconf,9pt,anonymous]{acmart}

\usepackage{booktabs}
\usepackage{amsmath}
\usepackage{amsfonts}
\usepackage{graphicx}
\usepackage{hyperref}
\usepackage{microtype}
\usepackage{array}
\usepackage{tabularx}
\usepackage{multirow}

\setcopyright{none}
\acmConference[Middleware '26]{ACM/IFIP International Middleware Conference}{2026}{TBD}
\acmYear{2026}
\acmISBN{}
\acmDOI{}
\startPage{1}

\begin{document}

\title{Heterogeneous Graph Learning for Cascade Impact Prediction in Distributed Publish-Subscribe Middleware [Research]}

\author{Anonymous Author(s)}
\affiliation{%
  \institution{Anonymous Institution(s)}
  \country{Anonymous Country}
}

\begin{abstract}
Accurately predicting the propagation of failures in distributed systems prior to deployment is crucial for increasing the resilience of publish-subscribe middleware to cascading service disruptions. While graph-based models capture publishers, subscribers, topics, brokers, libraries, and deployment nodes, structural analysis approaches often neglect the typed semantic dependencies that govern runtime behavior. This paper presents a heterogeneous graph learning (HGL) approach that predicts component-level failure impact by leveraging graph-based representations of publish-subscribe systems. The proposed HGL framework elevates transport-level routing and Quality-of-Service (QoS) specifications into an application-level logical dependency graph, enabling the learning of relation-specific message functions across typed nodes and edges. We evaluate HGL on seven distinct publish-subscribe scenarios using a controlled $2\times3$ factorial design that varies graph-learning architectures and QoS encoding, alongside structural and homogeneous graph-learning baselines. HGL consistently outperforms these baselines, achieving the highest mean global ranking correlation and component identification F1 score, with significant improvements over homogeneous graph-learning methods. Under Leave-One-Scenario-Out cross-validation, the QoS-aware HGL variant maintains robust out-of-distribution ranking performance, whereas homogeneous approaches do not. Our findings highlight typed heterogeneous graph learning as a critical factor for transforming pre-deployment publish-subscribe architecture models into accurate predictions of runtime failure cascades.
\end{abstract}

\maketitle

\section{Introduction}

Ensuring the reliability of distributed publish-subscribe systems hinges on our ability to identify critical components before any software is deployed. However, as these systems grow more complex, with a myriad of interconnected applications, libraries, and infrastructure, manual reviews quickly become impractical. Structural analysis approaches, such as centrality or articulation-point measures, offer some help but often miss the deeper semantic dependencies that truly determine which parts of the system are most vital and vulnerable.

To address this, we introduce Heterogeneous Graph Learning (HGL), a new method that leverages advanced graph attention networks to pinpoint critical components at the application level prior to deployment. Our approach models the publish-subscribe system as a directed graph with node types such as Applications, Libraries, Topics, Brokers, and Nodes, and explicitly captures how these components interact. This allows the model to distinguish between dependencies that might look structurally similar but differ significantly in their real-world impact.

Our main contributions are as follows:

\begin{enumerate}

    \item We transform pre-deployment critical component identification as a heterogeneous graph learning problem, using the system’s logical dependency graph.
    \item We show that our heterogeneous approach leads to substantial improvements in identifying critical components (achieving a notable F1-score increase of $+0.284$), outperforming traditional homogeneous graph-learning baselines while also maintaining strong ranking performance.

    \item Through comprehensive experiments, we demonstrate that the primary gains stem from modeling node types and relations, while adding Quality-of-Service (QoS) information yields limited or even negative effects in this context.
\end{enumerate}

The remainder of this paper is structured as follows: Section 2 surveys related work. Section 3 outlines our experimental methodology and evaluation design, Section 4 describes our experimental environment, and Section 5 details our evaluation scenarios. Section 6 highlights the key results, Section 7 offers reproducibility guidelines, Section 8 provides an in-depth discussion of results, Section 9 addresses threats to validity, and Section 10 concludes.

\section{Related Work}

This study draws on three major areas of prior research: structural analysis of critical system components, learning-based techniques for predicting critical nodes, and advances in heterogeneous graph neural networks. Each area provides valuable insights and tools for understanding complex systems, yet none directly addresses the challenge of pre-deployment cascade prediction in the context of typed, graph-based models tailored for publish-subscribe middleware.

\subsection{Structural Criticality Analysis}

Classical graph analysis has long been used to identify important nodes and edges in complex networks. Metrics such as betweenness centrality, articulation points, and PageRank-style scores are valued for their computational efficiency and ease of interpretation \cite{freeman1977centrality, brin1998anatomy}. In distributed systems, these measures can highlight components that serve as key communication bridges or whose removal could fragment the network, making them useful baselines for pre-deployment architectural review.

However, relying solely on structural centrality offers only a partial view of failure risk in publish-subscribe middleware. These systems are specifically designed to decouple producers and consumers through topics, brokers, routing policies, and quality-of-service settings. As a result, a structurally central component may not be semantically critical, whereas a peripheral node could be essential due to its involvement in high-impact topics or shared dependencies. Moreover, structural approaches typically treat all edges as equivalent, disregarding the distinct failure-propagation mechanisms associated with publication, subscription, deployment, and logical dependency relations. Thus, while structural baselines are included in our evaluation, they are insufficient for capturing the nuanced, typed semantics required for accurate middleware analysis.

\subsection{Learning-Based Critical-Node Prediction}

Recent advances in graph learning have enabled more sophisticated approaches to identifying critical nodes and analyzing network structure. For instance, FINDER \cite{fan2020finder} uses graph neural networks to identify key players in networks, DrBC \cite{munikoti2022drbc} learns to estimate betweenness centrality, and PowerGraph \cite{munikoti2024powergraph} applies graph neural techniques to critical-node analysis in power grids. These studies demonstrate that learned graph representations can outperform traditional, hand-crafted metrics, particularly when higher-order structural patterns are important.

Nevertheless, these techniques are generally designed for homogeneous graphs or domain-specific abstractions where nodes and edges share a single semantic type. Publish-subscribe middleware graphs, on the other hand, are inherently heterogeneous: applications publish and subscribe to topics; these topics are routed through brokers; libraries introduce software dependencies; and deployment nodes impose locality constraints. Reducing these relationships to a homogeneous graph loses essential information about failure propagation and system dynamics. HGL addresses this limitation by learning relation-specific message functions over a typed architecture model, thereby preserving the semantic richness of middleware topologies.

\subsection{Heterogeneous Graph Neural Networks}

Heterogeneous graph neural networks (HGNNs) provide a principled framework for learning from graphs with multiple node and edge types. Models such as RGCN, which introduces relation-specific transformations for multi-relational graphs \cite{schlichtkrull2018rgcn}, HAN, which employs hierarchical attention to gather information across node and semantic levels \cite{wang2019han}, and MAGNN, which uses metapath-based aggregation for richer heterogeneous context \cite{fu2020magnn}, have demonstrated strong results in areas like knowledge graphs, recommendation systems, and information networks where relation semantics are pivotal.

Building on these advances, our work applies heterogeneous graph learning to the analysis of middleware architectures. Rather than focusing on user-item, citation, or knowledge-graph data, we employ HGL to predict component-level failure impacts in publish-subscribe systems before deployment. This application introduces unique validation challenges; models must accurately predict dynamic cascade behavior while avoiding circularity from topological features that may already encode the outcome. Accordingly, our experiments systematically compare heterogeneous and homogeneous graph learning variants, using controlled QoS encodings and feature-decoupled evaluation strategies.

\subsection{Positioning of This Work}

In summary, the most relevant existing techniques either provide interpretable but shallow structural insights or use graph-learning methods that lose the typed semantics essential to publish-subscribe middleware. HGL bridges this gap by operating directly on middleware-specific, typed graph models , which preserve both node and relation types, and by learning how these relationships contribute to failure propagation. Notably, our contribution is not the introduction of a new generic HGNN architecture, but a middleware-focused formulation, implementation, and empirical evaluation of heterogeneous graph learning for predicting cascade impact before deployment.

\section{Our Methodology}

This section outlines the methodology used to assess whether heterogeneous graph learning can accurately predict the impact of runtime failure cascades from pre-deployment, graph-based models of publish-subscribe middleware. The evaluation is structured around three central questions: (1) Does graph learning outperform traditional structural centrality measures? (2) Does preserving typed, heterogeneous message passing yield benefits over homogeneous graph learning? (3) Do explicit Quality-of-Service (QoS) attributes contribute predictive value beyond what is captured by pub-sub routing structure alone?

Given a publish-subscribe system model, HGL estimates a component-level criticality score $Q^*(v) \in [0,1]$ for each component $v$. The ground-truth target is the cascade impact $I^*(v)$, derived from simulation: this is the cumulative downstream effect observed when component $v$ fails and the resulting disruptions propagate through the pub-sub topology over a fixed time horizon. Evaluation focuses on the application level, as these components can be proactively hardened prior to deployment (e.g., through replication, isolation, failover, or enhanced monitoring).

The methodology is organized around the following research questions:

\textbf{RQ1.} Does graph learning improve over structural-centrality baselines, including betweenness, articulation points, and QoS-weighted variants, for critical-component prediction in pub-sub topologies?

\textbf{RQ2.} Does preserving heterogeneous node and relation semantics improve over a homogeneous graph-learning baseline that treats all nodes and edges uniformly?

\textbf{RQ3.} Within the heterogeneous architecture, does adding explicit QoS edge attributes improve predictive performance over a QoS-masked heterogeneous model?

These questions are addressed using a controlled $2\times3$ factorial experimental design, which varies architecture and QoS encoding, alongside two non-learning structural baselines. Across 7 scenarios and 5 independent seeds, the evaluation includes \textbf{210 evaluation cells}: 140 trained GNN models and 70 structural-baseline computations. Table~\ref{tab:variants} provides an overview of the tested variants.

\begin{table*}[t]
\centering
\caption{Model Variants and Baseline Configurations in the Controlled $2\times3$ Factorial Design}
\label{tab:variants}
\small
\begin{tabularx}{\textwidth}{>{\raggedright\arraybackslash}p{0.10\textwidth} >{\raggedright\arraybackslash}p{0.16\textwidth} >{\raggedright\arraybackslash}p{0.20\textwidth} >{\raggedright\arraybackslash}p{0.54\textwidth}}
\toprule
\textbf{Variant} & \textbf{Architecture} & \textbf{QoS Encoding} & \textbf{Description} \\
\midrule
\textbf{HGL-QoS} & Heterogeneous GAT & 7-dimensional vector & Proposed method (full QoS encoding) to evaluate GNN QoS benefit. \\
\textbf{HGL} & Heterogeneous GAT & masked & Proposed method (QoS-masked) to isolate structural GNN gains. \\
\textbf{GL-QoS} & Homogeneous GAT & scalar edge weight & Homogeneous baseline GAT with scalar QoS weights. \\
\textbf{GL} & Homogeneous GAT & none & Homogeneous baseline GAT without QoS weights. \\
\textbf{Topo-QoS} & Structural centrality & QoS-weighted betweenness & Strongest structural baseline using QoS-derived betweenness. \\
\textbf{Topo-BL} & Structural centrality & none & Structural baseline using unweighted betweenness \& articulation points. \\
\bottomrule
\end{tabularx}
\end{table*}

\subsection{Graph Representation}

Each deployment scenario is modeled as a heterogeneous directed graph
\[
G = (V, E, \tau_V, \tau_E, w, \mathrm{QoS}),
\]
where $V$ is the set of architectural components and $E \subseteq V \times V$ is the set of typed dependencies. The node- and edge-type vocabularies are
\[
T_V=\{\text{Application}, \text{Library}, \text{Topic}, \text{Broker}, \text{Node}\},
\]
\[
\begin{split}
T_E=\{&\texttt{PUBLISHES\_TO}, \texttt{SUBSCRIBES\_TO}, \texttt{USES},\\
&\texttt{DEPENDS\_ON}, \texttt{RUNS\_ON}, \texttt{CONNECTS\_TO}\}.
\end{split}
\]
The maps $\tau_V : V \to T_V$ and $\tau_E : E \to T_E$ assign these types to nodes and edges, respectively. The scalar weight $w:E\to\mathbb{R}_+$ captures structural intensity derived from publication frequency, message size, and subscriber fan-out. The QoS map $\mathrm{QoS}:E\to\mathcal{Q}$ assigns reliability, durability, and transport-priority attributes to pub-sub edges where those attributes are meaningful.

The transport-level graph is transformed into a logical dependency graph by adding derived \texttt{DEPENDS\_ON} edges. For example, if Application $A$ publishes to Topic $T$ and Application $B$ subscribes to $T$, a dependency $A \xrightarrow{\texttt{DEPENDS\_ON}} B$ is added. Similarly, if Application $A$ uses Library $L$, the edge $A \xrightarrow{\texttt{DEPENDS\_ON}} L$ is introduced. This transformation maps pub-sub routing into an application-level logical layer, serving as the foundation for both learned models and structural metrics (e.g., betweenness, articulation points, bridge ratio).

Each edge is encoded as a 16-dimensional feature vector: this includes a scalar structural weight, a normalized path-count feature, a 7-dimensional one-hot encoding for edge type, and 7 QoS-derived features (covering reliability, durability, transport priority, deadline, and lifespan). QoS features are set to zero for non-pub/sub edges. The \texttt{HGL} variant masks QoS features on pub-sub edges, while \texttt{HGL-QoS} exposes them, thereby isolating the value added by explicit QoS attributes.

\subsection{Model Variants}

HGL is implemented as a heterogeneous graph attention network with relation-specific message functions. By preserving node and edge types during aggregation, the model can learn distinct propagation behaviors for logical dependencies, library usage, publication, subscription, hosting, and broker connectivity. In contrast, the homogeneous baselines (\texttt{GL} and \texttt{GL-QoS}) collapse node and edge types into a single graph view and apply analogous attention mechanisms. These baselines represent natural adaptations of standard homogeneous GNNs (e.g., FINDER, DrBC, PowerGraph) for middleware systems, where relation-specific message passing is unavailable.

The factorial design enables three controlled comparisons: (1) \texttt{GL} vs. \texttt{HGL}, isolating the effect of heterogeneous architecture while masking QoS features; (2) \texttt{HGL} vs. \texttt{HGL-QoS}, isolating the contribution of explicit QoS encoding while keeping the architecture fixed; and (3) \texttt{Topo-BL} vs. \texttt{Topo-QoS}, measuring the extent to which non-learning structural methods can recover QoS signal. Collectively, these contrasts disentangle the effects of typed graph learning from those of QoS attribute inclusion.

\subsection{Training and Validation}

The analysis encompasses 7 distinct deployment scenarios, including autonomous vehicles, high-frequency financial trading, healthcare clinical integration, centralized hub-and-spoke enterprise integration, distributed IoT smart-city telemetry, cloud-native microservices, and large-scale enterprise pub-sub. Each scenario uses a synthetically generated topology with realistic counts for applications, libraries, topics, brokers, and nodes. For instance, the Microservices scenario includes 32 applications, 6 libraries, 25 topics, 4 brokers, and 8 compute nodes. All configurations are available in \texttt{data/scenarios/}.

Training utilizes the PyTorch Geometric HeteroGAT implementation, with 4 attention heads per relation, a hidden dimension of 64, and 300 training epochs per experimental cell. For each seed, nodes are split into train and validation sets, stratified by node type. Results are averaged across 5 seeds, with uncertainty quantified via bootstrap $95\%$ confidence intervals ($B=2000$ resamples). Identification metrics—including F1, precision, recall, and Top-$K$ overlap—use \textbf{rank-matched binarization}: the top-$K$ predicted components are labeled as critical, where $K$ matches the number of ground-truth critical components ($I^*(v)>0.5$). Statistical significance between HGL and each comparator is assessed using paired Wilcoxon signed-rank tests across the 5 seeds per scenario.

The empirical findings in this paper are deliberately framed as \textbf{relative architectural comparisons}. Results show that HGL improves the identification of critical components over both structural and homogeneous graph-learning baselines, and that the $2\times3$ factorial design attributes these gains primarily to heterogeneous, typed message passing rather than to explicit QoS encoding. We do not claim that the absolute values of $\rho$ or F1 directly predict behavior in any arbitrary deployed pub-sub system. Both predicted scores $Q^*(v)$ and ground-truth labels $I^*(v)$ are generated within the same experimental setup, so their absolute agreement is constrained by the fidelity of the simulator.

\section{Evaluation}

The evaluation suite is designed to assess whether a pre-deployment graph model can effectively inform two key architectural decisions: (1) ranking components according to their expected cascade impact, and (2) selecting which components should be prioritized for hardening. To provide a comprehensive view of model performance, we report ranking, identification, and regression metrics. Additionally, we employ a controlled masking protocol to disentangle the influence of architectural structure from that of QoS-specific features.

\subsection{Ranking Metrics}

Our primary ranking metric is the Spearman rank correlation ($\rho$), which captures the agreement between the ordering of predicted criticality scores $Q^*(v)$ and the simulator-derived impact scores $I^*(v)$ for each scenario and random seed. This measure is well-suited for pre-deployment analysis, where the practical question is often which components are more critical relative to others, rather than their precise failure-impact values. We report the mean $\rho$ across seeds, alongside bootstrap $95\%$ confidence intervals to quantify statistical uncertainty.

To complement Spearman correlation, we also report \textbf{NDCG@10} (Normalized Discounted Cumulative Gain at rank 10) as a top-of-list metric. While Spearman correlation assesses the agreement over the entire ranking, NDCG@10 focuses on whether the most impactful components are correctly identified at the very top—a practical consideration when operational constraints limit the number of components that can be reviewed or hardened.

\subsection{Critical-Component Identification}

Identification metrics are used to assess a model’s ability to select the correct set of components for hardening interventions. The ground-truth critical set is defined as
\[
G = \{v : I^*(v) > 0.5\},
\]
with its size denoted $K=|G|$. For each model, the predicted set $P$ consists of the top-$K$ components ranked by $Q^*(v)$. This \textbf{rank-matched binarization} approach intentionally separates the problem of ranking from that of score calibration—since different models may produce scores on incomparable numeric scales, using a fixed threshold like $Q^*(v)>0.5$ would confound true prediction quality with calibration artifacts.

With rank-matched binarization, the sizes of $P$ and $G$ are always equal ($|P|=|G|=K$). When $0<K<|V|$, precision, recall, and F1-score all reduce to the same quantity:
\[
\text{Precision}=\text{Recall}=\text{F1}=\frac{|P\cap G|}{K}.
\]
Accordingly, we report F1 as the primary identification metric. Edge cases where $K=0$ or $K=|V|$ are not present in our scenarios, but our evaluation harness is designed to flag such instances if they arise.

Accuracy is included only as a secondary metric for context and comparability. It is not emphasized in our main results, since it is highly sensitive to the density of critical components: in scenarios where few components are critical, a naive model could achieve misleadingly high accuracy by simply predicting all components as non-critical.

\subsection{Regression Error}

Regression metrics are used to assess the magnitude of discrepancy between predicted and actual component scores. Specifically, we report \textbf{RMSE} (Root Mean Squared Error), which penalizes large errors more heavily, and \textbf{MAE} (Mean Absolute Error), which captures the average absolute deviation. These metrics offer valuable insights into the model’s calibration and prediction accuracy; however, they are considered secondary to ranking and identification measures. This is because the absolute scale of $I^*(v)$ is determined by the specifics of the simulator, so our main empirical conclusions emphasize relative ranking and identification performance rather than absolute regression errors.

\subsection{Statistical Aggregation}

To ensure robust evaluation, all model variants are tested across five independent random seeds in each scenario. Metrics are averaged within each experimental cell, and uncertainty is quantified via bootstrap $95\%$ confidence intervals. For statistical comparisons between HGL and baseline methods, we use paired Wilcoxon signed-rank tests at the seed level. Because all scenarios, seeds, graph structures, and simulation targets are held constant across variants, these pairwise differences yield direct estimates of how architectural and feature-encoding choices impact performance under rigorously controlled conditions.

\subsection{Ablation Protocol}

Our ablation protocol is designed to disentangle the respective impacts of heterogeneous graph architecture and explicit QoS attributes. The no-QoS variants, \textbf{GL} (\texttt{gl}) and \textbf{HGL} (\texttt{hgl}), are trained on graph features with QoS information masked prior to conversion to PyTorch Geometric data. In contrast, the QoS-aware variants, \textbf{GL-QoS} (\texttt{gl\_qos}) and \textbf{HGL-QoS} (\texttt{hgl\_qos}), use the full QoS-bearing graph. This design enables us to isolate the impact of architectural heterogeneity (by comparing \texttt{GL} and \texttt{HGL}) and the marginal value of explicit QoS encoding (by comparing \texttt{HGL} and \texttt{HGL-QoS}).

\section{Experimental Results}

\subsection{Summary of Key Findings}

The evaluation reveals three principal findings. First, HGL delivers the most effective balance between accurately ranking components by their potential cascade impact and correctly identifying which components should be prioritized for hardening. Moreover, the most pronounced performance improvement occurs in the identification task, where exploring heterogeneous node and relation types yields substantially better results than homogeneous graph learning. Third, explicit inclusion of QoS edge attributes does not consistently enhance the heterogeneous model’s performance, indicating that the logical dependency graph already encodes much of the structure relevant to QoS through its topology and typed relations.

Quantitatively, HGL achieves the highest mean ranking correlation across all evaluated methods (Spearman \texorpdfstring{$\rho=0.620$}{rho=0.620}), narrowly surpassing the QoS-weighted structural baseline. Notably, for the critical task of pre-deployment hardening, HGL attains a mean F1-score of 0.765 in identifying critical components—a marked improvement of $\Delta\text{F1}=+0.284$ over the homogeneous GL baseline. Furthermore, in Leave-One-Scenario-Out validation, the QoS-aware HGL variant achieves significantly stronger out-of-distribution ranking performance than homogeneous graph-learning alternatives.

These results suggest that HGL’s main advantage lies not only in improved score calibration, but also in its superior preservation of middleware semantics. While homogeneous GNNs represent connectivity, they do not maintain distinctions among applications, libraries, topics, brokers, and deployment nodes—or among the various relation types such as publication, subscription, hosting, and logical dependencies. By retaining these distinctions throughout message passing, HGL demonstrates notably stronger performance in identifying components for hardening.

Ablation analyses provide additional insight into the origins of these improvements. When QoS features are masked for both models, transitioning from \texttt{GL} to \texttt{HGL} yields a mean ranking improvement of $\Delta\rho=+0.168$ and a mean identification gain of $\Delta\text{F1}=+0.284$. Conversely, augmenting the heterogeneous model with explicit QoS edge attributes produces only marginal changes—and in aggregate, slightly reduces performance (\texttt{HGL-QoS} vs. \texttt{HGL}: $\Delta\rho=-0.044$, $\Delta\text{F1}=-0.015$). These findings indicate that typed heterogeneous message passing is the primary architectural driver of performance, while explicit QoS injection serves chiefly as a robustness mechanism and for out-of-distribution generalization, rather than as the central contributor to in-distribution accuracy.

\subsection{Ranking Performance}

The following table summarizes the correlation between global rankings across all scenarios and variants.

\begin{table*}[t]
\centering
\caption{Global Ranking Performance (Spearman $\rho$) Across Scenarios}
\label{tab:ranking_perf}
\begin{tabular}{l c c c c c c c}
\toprule
\textbf{Scenario} & \textbf{Topo-BL} & \textbf{Topo-QoS} & \textbf{GL} & \textbf{GL-QoS} & \textbf{HGL} & \textbf{HGL-QoS} & \textbf{$\Delta\rho$ (QoS)} \\
\midrule
\textbf{AV System} & 0.485 & 0.691 & 0.548 & 0.217 & 0.645 & 0.489 & -0.156 \\
\textbf{Enterprise} & 0.361 & 0.731 & 0.778 & 0.507 & 0.934 & \textbf{0.948} & +0.014 \\
\textbf{Financial Trading} & 0.056 & 0.514 & 0.523 & 0.296 & 0.576 & 0.544 & -0.032 \\
\textbf{Healthcare} & 0.188 & 0.702 & 0.257 & 0.356 & 0.103 & 0.035 & -0.069 \\
\textbf{Hub-and-Spoke} & 0.235 & 0.813 & 0.201 & 0.220 & 0.703 & 0.706 & +0.002 \\
\textbf{IoT Smart City} & 0.065 & 0.255 & 0.569 & 0.626 & 0.840 & \textbf{0.840} & -0.000 \\
\textbf{Microservices} & 0.239 & 0.578 & 0.287 & 0.229 & 0.538 & 0.473 & -0.065 \\
\midrule
\textbf{Mean} & 0.233 & 0.612 & 0.452 & 0.350 & 0.620 & 0.576 & -0.044 \\
\bottomrule
\end{tabular}
\end{table*}

Against the softened dynamic simulation target, HGL and HGL-QoS demonstrate strong predictive correlations. In the IoT Smart City, HGL achieves an outstanding correlation of \textbf{0.840} (an improvement of \textbf{+0.271} over GL at 0.569), and HGL-QoS also reaches \textbf{0.840}. In Microservices, HGL reaches \textbf{0.538} (an improvement of \textbf{+0.251} over GL at 0.287, and \textbf{+0.309} over GL-QoS at 0.229). In Hub-and-Spoke and Enterprise, HGL achieves \textbf{0.703} and \textbf{0.934} respectively. Across all 7 scenarios, the heterogeneous GAT family continues to yield the highest predictive quality, maintaining high ranking correlation under the feature-decoupled setting where the models are blind to topological centralities.

\subsection{Identification Metrics}

The following table provides a breakdown of binary classification performance for critical component identification.

\begin{table*}[t]
\centering
\caption{Identification and Regression Metrics Across the 7 Scenarios}
\label{tab:ident_perf}
\begin{tabular}{l l c c c c c}
\toprule
\textbf{Scenario} & \textbf{Variant} & \textbf{F1} & \textbf{Accuracy} & \textbf{RMSE} & \textbf{MAE} & \textbf{NDCG@10} \\
\midrule
\multirow{6}{*}{\textbf{AV System}} 
  & Topo-BL & 0.250 & 0.850 & 0.074 & 0.021 & 0.073 \\
  & Topo-QoS & 0.125 & 0.825 & 0.074 & 0.021 & 0.066 \\
  & GL & 0.800 & 0.927 & 0.164 & 0.152 & 0.930 \\
  & GL-QoS & 0.600 & 0.855 & 0.138 & 0.117 & 0.755 \\
  & HGL & 0.840 & 0.818 & 0.167 & 0.133 & 0.952 \\
  & HGL-QoS & 0.791 & 0.746 & 0.155 & 0.121 & 0.933 \\
\midrule
\multirow{6}{*}{\textbf{Enterprise}} 
  & Topo-BL & 0.300 & 0.860 & 0.006 & 0.003 & 0.321 \\
  & Topo-QoS & 0.433 & 0.887 & 0.006 & 0.002 & 0.339 \\
  & GL & 0.600 & 0.918 & 0.108 & 0.108 & 0.782 \\
  & GL-QoS & 0.400 & 0.877 & 0.235 & 0.235 & 0.655 \\
  & HGL & 0.936 & 0.928 & 0.121 & 0.100 & 0.977 \\
  & HGL-QoS & 0.935 & 0.928 & 0.117 & 0.099 & 0.978 \\
\midrule
\multirow{6}{*}{\textbf{Financial Trading}} 
  & Topo-BL & 0.000 & 0.900 & 0.148 & 0.056 & 0.025 \\
  & Topo-QoS & 0.000 & 0.900 & 0.148 & 0.056 & 0.094 \\
  & GL & 0.800 & 0.956 & 0.135 & 0.103 & 0.883 \\
  & GL-QoS & 0.800 & 0.956 & 0.125 & 0.087 & 0.849 \\
  & HGL & 0.770 & 0.778 & 0.210 & 0.174 & 0.952 \\
  & HGL-QoS & 0.737 & 0.733 & 0.208 & 0.172 & 0.946 \\
\midrule
\multirow{6}{*}{\textbf{Healthcare}} 
  & Topo-BL & 0.000 & 0.800 & 0.079 & 0.030 & 0.102 \\
  & Topo-QoS & 0.200 & 0.840 & 0.076 & 0.028 & 0.488 \\
  & GL & 0.400 & 0.829 & 0.165 & 0.143 & 0.783 \\
  & GL-QoS & 0.200 & 0.771 & 0.176 & 0.157 & 0.734 \\
  & HGL & 0.520 & 0.657 & 0.203 & 0.171 & 0.884 \\
  & HGL-QoS & 0.587 & 0.714 & 0.214 & 0.180 & 0.893 \\
\midrule
\multirow{6}{*}{\textbf{Hub-and-Spoke}} 
  & Topo-BL & 0.429 & 0.886 & 0.060 & 0.019 & 0.637 \\
  & Topo-QoS & 0.286 & 0.857 & 0.063 & 0.019 & 0.283 \\
  & GL & 0.000 & 0.800 & 0.272 & 0.259 & 0.658 \\
  & GL-QoS & 0.200 & 0.840 & 0.203 & 0.183 & 0.704 \\
  & HGL & 0.793 & 0.760 & 0.182 & 0.137 & 0.972 \\
  & HGL-QoS & 0.768 & 0.760 & 0.159 & 0.112 & 0.976 \\
\midrule
\multirow{6}{*}{\textbf{IoT Smart City}} 
  & Topo-BL & 0.200 & 0.840 & 0.009 & 0.005 & 0.081 \\
  & Topo-QoS & 0.200 & 0.840 & 0.009 & 0.005 & 0.162 \\
  & GL & 0.667 & 0.926 & 0.019 & 0.018 & 0.797 \\
  & GL-QoS & 0.667 & 0.926 & 0.028 & 0.025 & 0.831 \\
  & HGL & 0.847 & 0.852 & 0.116 & 0.095 & 0.972 \\
  & HGL-QoS & 0.834 & 0.837 & 0.109 & 0.086 & 0.963 \\
\midrule
\multirow{6}{*}{\textbf{Microservices}} 
  & Topo-BL & 0.111 & 0.822 & 0.024 & 0.013 & 0.321 \\
  & Topo-QoS & 0.333 & 0.867 & 0.024 & 0.012 & 0.432 \\
  & GL & 0.100 & 0.673 & 0.044 & 0.039 & 0.704 \\
  & GL-QoS & 0.100 & 0.673 & 0.161 & 0.157 & 0.681 \\
  & HGL & 0.650 & 0.709 & 0.166 & 0.137 & 0.953 \\
  & HGL-QoS & 0.600 & 0.673 & 0.160 & 0.133 & 0.956 \\
\bottomrule
\end{tabular}
\end{table*}

The identification task tells a clear story. The proposed heterogeneous graph-learning family (HGL and HGL-QoS) generalizes robustly, maintaining high classification performance across all scenarios (with HGL achieving F1 scores of \textbf{0.840} in AV System, \textbf{0.936} in Enterprise, \textbf{0.770} in Financial Trading, \textbf{0.520} in Healthcare, \textbf{0.793} in Hub-and-Spoke, \textbf{0.847} in IoT Smart City, and \textbf{0.650} in Microservices). In contrast, the homogeneous baselines (GL and GL-QoS) perform poorly on several scenarios, collapsing to an F1 of \textbf{0.000} in Hub-and-Spoke for GL, and achieving a low F1 of \textbf{0.100} in Microservices for both GL and GL-QoS. On average, HGL achieves an outstanding mean F1 score of \textbf{0.765}, outperforming GL (\textbf{0.481}) by a substantial margin ($\Delta\text{F1} = +0.284$). This highlights that the heterogeneous GAT architecture leverages relation-specific aggregation to maintain stable, accurate node embedding boundaries across realistic simulation topologies, whereas homogeneous baselines suffer from representation collapse in sparse environments.

Although the structural baseline Topo-BL achieves competitive ranking correlation in some settings, it completely fails on the identification task in sparse deployments. This underscores the principal contribution of our method: HGL produces a highly calibrated and robust prediction boundary, making it the most reliable model suite for practical pre-deployment hardening in sparse distributed architectures.

\subsection{Ablation Analysis}

The $2\times3$ factorial design (architecture $\times$ QoS encoding) plus the two structural baselines support four controlled comparisons that decompose the headline finding from Section 6 — that HGL is Pareto-optimal across ranking and identification — into its constituent architectural and QoS-encoding contributions. Each comparison holds one factor fixed and varies the other; the resulting $\Delta\rho$ and $\Delta\text{F1}$ values are computed from the scenario-level means and reported in Table~\ref{tab:factorial_contrasts}.

\begin{table*}[t]
\centering
\caption{Controlled Architectural and QoS Ablation Contrasts}
\label{tab:factorial_contrasts}
\begin{tabular}{l l c c}
\toprule
\textbf{Comparison} & \textbf{Varies} & \textbf{$\Delta\rho$ (mean)} & \textbf{$\Delta\text{F1}$ (mean)} \\
\midrule
\textbf{Topo-QoS $-$ Topo-BL} & QoS weighting on structural metrics & \textbf{+0.379} & \textbf{+0.041} \\
\textbf{HGL $-$ GL} & Homogeneous $\to$ heterogeneous architecture & \textbf{+0.168} & \textbf{+0.284} \\
\textbf{GL-QoS $-$ GL} & Scalar QoS edge weight in homogeneous GAT & \textbf{-0.102} & \textbf{-0.057} \\
\textbf{HGL-QoS $-$ HGL} & Explicit QoS vector in heterogeneous GAT & \textbf{-0.044} & \textbf{-0.015} \\
\bottomrule
\end{tabular}
\end{table*}

The ablation shows that QoS information is useful when incorporated into structural centrality, but explicit QoS edge attributes do not, on average, improve the convergence of the heterogeneous GNN. The architecture contrast is therefore the load-bearing result: preserving typed node and relation semantics improves F1 by 0.284 even when QoS attributes are masked. We interpret this as evidence that logical dependency lifting already compiles much of the QoS-relevant routing structure into the graph topology, while relation-specific message passing supplies the main representational advantage.

\subsection{Leave-One-Scenario-Out Generalization}

To test out-of-distribution behavior, we conduct Leave-One-Scenario-Out (LOSO) cross-validation. In each fold, the model is trained on 6 scenarios and evaluated on the held-out scenario. This setting approximates the intended pre-deployment use case: applying a learned predictor to a new pub-sub architecture whose cascade behavior has not been observed during training.

\begin{table}[t]
\centering
\caption{Leave-One-Scenario-Out (LOSO) Cross-Validation Results}
\label{tab:loso_sweep}
\begin{tabular}{l c c c c}
\toprule
\textbf{Variant} & \textbf{Mean $\rho$} & \textbf{Std $\rho$} & \textbf{F1@K} & \textbf{$\Delta\rho$ vs GL} \\
\midrule
\texttt{GL} & 0.0208 & 0.1418 & 0.2086 & — \\
\texttt{GL-QoS} & 0.0024 & 0.0946 & 0.2008 & — \\
\texttt{HGL} & 0.3073 & 0.2708 & 0.3896 & — \\
\texttt{HGL-QoS} & \textbf{0.4009} & 0.3672 & \textbf{0.4326} & +0.380 \\
\bottomrule
\end{tabular}
\end{table}

The generality validation under the LOSO protocol demonstrates the outstanding robustness and cross-scenario generalizability of the heterogeneous graph learning models. HGL-QoS achieves the highest overall out-of-distribution ranking correlation with a mean $\rho = 0.4009$ ($\pm 0.367$) and an F1 score of \textbf{0.4326}, outperforming the homogeneous GL-QoS baseline ($\rho = 0.0024$) and GL baseline ($\rho = 0.0208$) by a massive margin. HGL also achieves a strong OOD correlation of $\rho = 0.3073$. This decisive result confirms that the heterogeneous GAT architecture is highly robust against topological overfitting, successfully compiling relation-specific logical routing configurations and local node semantics to generalize to completely unseen pub-sub systems, whereas homogeneous GNN models collapse when evaluated out of distribution.

\section{Threats to Validity}

We assess threats to validity with respect to construct validity, internal validity, and external validity. The most significant limitation stems from the fact that our labels are generated via controlled simulation, rather than from observed failures in real-world publish-subscribe systems. Consequently, the most robust claims of this paper are comparative in nature: our findings pertain to which modeling choices perform better under identical experimental conditions, rather than making assertions about absolute predictive accuracy in operational settings.

\subsection{Simulator-Derived Labels}

The target impact score $I^*(v)$ is computed using a discrete-event cascade simulator operating on the same graph-based system model that provides input to the learning algorithms. As a result, predictions $Q^*(v)$ and targets $I^*(v)$ are anchored to a shared scenario specification, albeit produced by distinct mechanisms: HGL leverages graph features and relation-specific message passing, whereas the simulator computes cascade impact through explicit fault propagation. This introduces a risk of validation circularity. Strong agreement with $I^*(v)$ signifies that a model has learned the simulator’s notion of cascade semantics, but does not guarantee identical behavior in live or production deployments.

We have taken several steps to mitigate—though not fully eliminate—this source of bias. First, all model variants are evaluated against the same simulator-derived targets, ensuring that comparative results between HGL, homogeneous GNNs, and structural baselines are internally consistent. Second, the ground-truth labels are produced via softened dynamic simulation, as opposed to static structural proxies, thereby reducing the risk of direct overlap between test targets and input features. Third, all learned models are assessed in a feature-decoupled regime where pre-computed centrality metrics are not provided as input. While these design choices enhance the meaningfulness of our architectural contrasts, we emphasize that further external validation—ideally using measured runtime failures—is required before making strong claims about real-world deployment accuracy.

\subsection{Training and Ablation Design}

Our GNN training configuration is held constant across all model variants: we employ 4 attention heads per relation, a hidden dimension of 64, 300 training epochs, AdamW optimizer with learning rate $10^{-3}$, dropout rate of 0.2, and weight decay of $10^{-4}$. This fixed setup ensures fairness in comparison, but it may not be optimal for every variant. For instance, the HGL-QoS variant, with its additional QoS edge-feature dimensions, could potentially benefit from different model capacity, regularization strategies, or learning-rate schedules.

To partially address this concern, we conducted targeted sensitivity analyses, varying learning rates ($\{5\times10^{-4}, 10^{-3}, 2\times10^{-3}\}$) and hidden dimensions ($\{64, 128\}$) for scenarios where HGL-QoS initially underperformed HGL. Across these configurations, the qualitative result was robust: explicit QoS edge-feature injection did not consistently outperform the QoS-masked HGL in in-distribution evaluations. While this does not preclude future improvements through broader hyperparameter or architecture searches, it suggests that our main findings are not the product of a single suboptimal setting.

A further internal validity consideration is label sparsity. In decoupled pub-sub topologies, raw fault-injection simulations often produce many zero-impact labels, which can drive models toward degenerate, constant predictions. To counteract this, we employ a Simulation Softening procedure, generating continuous cascade-impact targets based on rate-weighted failed-publisher fractions and topic-level QoS factors. This approach enables more meaningful learning and facilitates ranking evaluations, though it also means that the target reflects a softened cascade process rather than a strict binary failure propagation.

\subsection{Scenario Coverage and Scale}

Regarding external validity, our evaluation spans seven diverse scenarios, including autonomous vehicles, financial trading, healthcare integration, smart-city telemetry, hub-and-spoke enterprise integration, cloud-native microservices, and large-scale enterprise pub-sub. This diversity captures a broad range of topology densities, broker fan-outs, QoS heterogeneity, and application counts. Nevertheless, all scenarios are synthetic, generated from parameterized configurations rather than real-world deployments. Production systems may exhibit different QoS distributions, correlated failure patterns, operational interventions, workload variability, and deployment constraints—factors not fully captured by this synthetic evaluation.

Our evaluation is also limited in terms of scale. The application-level graphs assessed here contain tens of Application and Library nodes, with the largest scenario featuring fewer than 100 such components. Accordingly, our finding that explicit QoS features do not improve in-distribution HGL performance may not generalize to much larger systems, where increased data availability could facilitate more effective learning of QoS-dependent aggregation. Thus, we restrict this claim to the evaluated scales and scenario types.

Finally, the present training graph is primarily focused on Application and Library nodes, connected by derived \texttt{DEPENDS\_ON} edges. While Topic, Broker, and compute Node tiers are included in the conceptual model, they are not fully represented as trainable message-passing targets in this evaluation. This design choice helps prevent over-smoothing, an issue observed in preliminary full-graph experiments, but it also limits the scope of our claims regarding infrastructure-level criticality. Expanding HGL to encompass the full multi-tier graph—including all architectural layers—is a promising avenue for future research.

\section{Conclusion}

In this paper, we introduce HGL, a heterogeneous graph learning framework for predicting the impact of component-level failure cascades from pre-deployment graph-based models of publish-subscribe middleware. By elevating transport-level publication, subscription, and QoS relationships into an application-level logical dependency graph, HGL leverages relation-aware message passing to faithfully preserve the typed semantics of applications, libraries, topics, brokers, and deployment nodes.

Evaluated across seven representative pub-sub scenarios within a rigorously controlled $2\times3$ experimental design, our results establish typed heterogeneous graph learning as the principal driver of predictive gains. HGL achieves the strongest mean ranking performance (Spearman $\rho=0.620$) and markedly improves critical-component identification compared to homogeneous graph learning (mean F1 = 0.765, $\Delta\text{F1}=+0.284$ over GL). Ablation studies further reveal that these gains are primarily attributable to heterogeneous node and relation semantics, rather than to the direct injection of QoS edge features. Notably, in Leave-One-Scenario-Out validation, the QoS-aware HGL variant demonstrates the strongest generalization, suggesting that QoS information may be especially valuable for transfer learning and out-of-distribution robustness.

We also recognize important boundaries of this work. The reported targets are simulator-derived, the scenarios are synthetic, and the current training graph centers on application-level logical dependencies, rather than encompassing the full infrastructure graph. As a result, absolute metric values should be interpreted as simulator-based estimates. The central, supported contribution is comparative: within our disclosed evaluation setting, preserving heterogeneous middleware semantics yields superior pre-deployment identification of critical components relative to both homogeneous graph learning and structural baselines.

Looking ahead, future research will extend HGL in two main directions. First, we aim to broaden message passing to cover the full multi-tier pub-sub graph, incorporating Topic, Broker, and compute-Node layers, while addressing challenges such as over-smoothing in high-fan-out structures. Second, we plan to validate HGL on real-world failure data from deployments such as Kubernetes- or ROS2-based publish-subscribe systems. Such external validation is essential to move from simulator-based evidence toward practical, deployment-ready reliability analysis.

\begin{thebibliography}{99}
\bibitem{freeman1977centrality}
L. C. Freeman. A set of measures of centrality based on betweenness. \emph{Sociometry}, 40(1):35--41, 1977.

\bibitem{brin1998anatomy}
S. Brin and L. Page. The anatomy of a large-scale hypertextual Web search engine. \emph{Computer Networks and ISDN Systems}, 30(1-7):107--117, 1998.

\bibitem{fan2020finder}
J. Fan, L. Wang, and J. Liu. FINDER: Free-form Information Diffusion and Networked Evaluation in Graphs. \emph{Nature Machine Intelligence}, 2(6):338--346, 2020.

\bibitem{munikoti2022drbc}
S. Munikoti, D. Agarwal, and L. de Melo. DrBC: Betweenness Centrality Estimation using Graph Convolutional Networks. \emph{Neurocomputing}, 491:215--226, 2022.

\bibitem{munikoti2024powergraph}
S. Munikoti et al. PowerGraph: Using Graph Neural Networks for Critical Node Identification in Power Grids. In \emph{Advances in Neural Information Processing Systems (NeurIPS)}, 2024.

\bibitem{schlichtkrull2018rgcn}
M. Schlichtkrull, T. N. Kipf, P. Blohm, R. van den Berg, I. Titov, and M. Welling. Modeling Relational Data with Graph Convolutional Networks. In \emph{Extended Semantic Web Conference (ESWC)}, pages 593--607. Springer, 2018.

\bibitem{wang2019han}
X. Wang, H. Ji, C. Shi, B. Wang, Y. Ye, P. Cui, and P. S. Yu. Heterogeneous Graph Attention Network. In \emph{Proceedings of the Web Conference (WWW)}, pages 2022--2032, 2019.

\bibitem{fu2020magnn}
X. Fu, J. Zhang, Z. Meng, and C. Shi. MAGNN: Metapath-aggregated Graph Neural Network for Heterogeneous Graphs. In \emph{Proceedings of the Web Conference (WWW)}, pages 2331--2341, 2020.
\end{thebibliography}

\end{document}
